% from main.tex
%%%%%%%%%%  scan idx 284  |  folio 271  %%%%%%%%%%
localement constante d'après (5.5 (ii)). L'assertion de constructibilité de (i) s'en déduit aussitôt. Les autres assertions découlent de (5.5 (ii) et (iii)) et (5.2 (iv)), sauf la réciproque de (ii), pour laquelle il suffit, après s'être ramené à \(E\) strictement parfait (après localisation et troncation), d'appliquer (EGA III 7.6.2 et 7.6.5) au foncteur \(T_{-p}=H^p(E\otimes\,\cdot\,)\).

\textbf{Proposition 5.7.} (EGA III 7.7.5 et 7.9.4). Soient \(f:X\to Y\) un morphisme propre de schémas quasi-compacts, \(E\in\Ob D^b(f)_{\parf}\) (4.2), et \(p\in\mathbb Z\). Alors, sous réserve de validité de (2.1), donc par exemple si \(f\) est projectif ou si \(Y\) est noethérien :
\begin{enumerate}[label=(\roman*)]
\item Pour \(y\in Y\), les \(H^i(Rf_*(E)\otimes_Y^{\mathrm L}k(y))\) sont des espaces vectoriels de rang fini sur \(k(y)\), nuls sauf un nombre fini, et la fonction
\[
 y\longmapsto\sum(-1)^i\operatorname{rg}_{k(y)}H^i(Rf_*(E)\otimes_Y^{\mathrm L}k(y))
\]
est localement constante sur \(Y\).
\item La fonction
\[
 y\longmapsto\operatorname{rg}_{k(y)}H^p(Rf_*(E)\otimes_Y^{\mathrm L}k(y))
\]
est constructible et semi-continue supérieurement sur \(Y\).
\item Supposons \(Y\) quasi-séparé. Alors, si le foncteur \(M\mapsto R^pf_*(E\otimes_Y^{\mathrm L}M)\) est exact sur la catégorie des \(\mathcal O_Y\)-Modules quasi-cohérents, \(R^pf_*(E)\) est un \(\mathcal O_Y\)-Module localement libre de type fini, et la fonction de (ii) est localement constante. Inversement, si \(Y\) est noethérien et réduit, et si la fonction de (ii) est localement constante, le foncteur \(M\mapsto R^pf_*(E\otimes_Y^{\mathrm L}M)\) est exact sur la catégorie des \(\mathcal O_Y\)-Modules quasi-cohérents.
\end{enumerate}

